

\section{Addestramento}

Per la fase di training andremo ad utilizzare dati generati da Quantum Espresso e un potenziale GAP. Partiamo esponendo in breve la struttura del processo che andremo ad utilizzare.

La funzione svolta dal potenziale GAP è quella di associare determinate energie e forze, fornite come detto da Qauntum Espresso, ad una specifica configurazione geometrica, questa descritta tramite SOAP. 

Questa associazione viene poi utilizzata per indovinare il valore di energia di un atomo dato il suo introno geometrico, in particolare, preso un atomo X ne viene calcolato il SOAP, questo, tramite una funzione di kernel, viene confrontato con gli ambienti della fase di training ed infine viene estratto un valore probabile per l'energia.
Data la struttura del metodo, un punto che vogliamo sottolineare è che l'accuratezza decresce velocemente fuori da scenari compatibili con i dati di training, è ovviamente possibile arginare se non risolvere il problema addizionando diverse tecnologie al metodo, queste non sono però esaminate in questa trattazione.
